\documentclass[interno]{gpt}

\titolo{Unione e analisi preliminare dei requisiti}
\data{2026-11-16}
\orario{18:00}{20:00}
\luogo{Discord}
\redattore{guglielmo}
\verificatore{biasini}
%\assenti{}

\componenti{biasini,cescobolla,corro,guglielmo,mazzurana,spagnoli}
\begin{document}

\section{Ordine del giorno}
\begin{enumerate}
  \item Unione dei requisiti individuati individualmente per sezioni.
  \item Revisione e discussione di gruppo dei requisiti.
  \item Preparazione dell'incontro con il proponente di mercoledì 18 novembre 2026.
  \item Organizzazione della retrospettiva dello sprint 1.
\end{enumerate}

\clearpage

\section{Svolgimento}

\subsection{Unione dei requisiti individuati per sezioni}
Nei giorni precedenti ogni componente ha analizzato in autonomia una sezione del capitolato e ha individuato i requisiti relativi.
Durante la riunione il gruppo ha unito questi contributi in un'unica lista, il più possibile dettagliata.
La lista ha carattere preliminare: l'\textit{Analisi dei requisiti} è ancora in fase di definizione.

\subsection{Revisione e discussione di gruppo dei requisiti}
Il gruppo ha rivisto e discusso insieme ogni requisito della lista.
Sono state eliminate le ridondanze, cioè i requisiti emersi in più sezioni, e le voci ritenute superflue.
Per gli altri requisiti il gruppo ha cercato di scendere il più possibile nel dettaglio, così da arrivare all'incontro con il proponente con domande precise.

\subsection{Preparazione dell'incontro con il proponente}
La lista preliminare sarà discussa con Eggon nell'incontro di mercoledì 18 novembre 2026 alle 17:00.
Solo dopo il confronto con il proponente il gruppo procederà alla redazione vera e propria dell'\textit{Analisi dei requisiti}, completa dei casi d'uso.
In questo modo si evita di formalizzare requisiti che il proponente potrebbe modificare o scartare.

\subsection{Organizzazione della retrospettiva dello sprint 1}
L'analisi preliminare dei requisiti ha richiesto più tempo del previsto e ha occupato l'intera riunione.
Il gruppo ha quindi deciso di rinviare la retrospettiva dello sprint 1 a una riunione straordinaria.
La riunione si terrà mercoledì 18 novembre 2026, subito dopo l'incontro con Eggon.

\section{Decisioni}
\begin{decisioni}
    \decisione{Adozione di un'unica lista preliminare dei requisiti, ottenuta dall'unione dei contributi individuali e ripulita da ridondanze e voci superflue.}
    \decisione{Redazione completa dell'\textit{Analisi dei requisiti}, con i casi d'uso, solo dopo l'incontro con il proponente di mercoledì 18 novembre 2026.}
    \decisione{Retrospettiva dello sprint 1 in una riunione straordinaria, mercoledì 18 novembre 2026 dopo l'incontro con Eggon.}
\end{decisioni}

\section{Azioni da intraprendere}
\begin{azioni}
  \azione{Stesura verbale interno della riunione odierna.}{guglielmo}{17/11/2026}
  \azione{Discussione della lista preliminare dei requisiti con il proponente.}{biasini,cescobolla,corro,guglielmo,mazzurana,spagnoli}{18/11/2026}
  \azione{Retrospettiva dello sprint 1 in riunione straordinaria.}{biasini,cescobolla,corro,guglielmo,mazzurana,spagnoli}{18/11/2026}
  \azione{Redazione dell'\textit{Analisi dei requisiti} con i casi d'uso, in base all'esito dell'incontro con il proponente.}{corro,spagnoli}{29/11/2026}
\end{azioni}

\section{Prossima riunione}
La prossima riunione interna è una riunione straordinaria, prevista per mercoledì 18 novembre 2026 subito dopo l'incontro con Eggon delle 17:00.
Lo scopo è svolgere la retrospettiva dello sprint 1.

\end{document}
